\begin{table}[htbp]
\centering
\caption{\\ Stock market responses to key legal events}
\label{tab:event_study}
\small
\setlength{\tabcolsep}{4pt}
\medskip
\begin{adjustbox}{max width=\textwidth}
\begin{tabular}{lcccccc}
\toprule
Event & Dates & Days & Market & No gold & Gold & Gold \\
 & & & & (equal-wt.) & (equal-wt.) & (exposure-wt.) \\
\midrule
Joint Resolution & May 26--June 6, 1933 & 9 & +12.1\% & +25.4\% & +30.9\% & +31.4\% \\
Supreme Court oral arguments & Jan.~8--10, 1935 & 3 & -1.3\% & -1.4\% & -1.3\% & -1.4\% \\
Supreme Court decision & Feb.~18, 1935 & 1 & +2.9\% & +3.2\% & +3.7\% & +3.5\% \\
\bottomrule
\end{tabular}
\end{adjustbox}
\medskip
\begin{minipage}{0.95\textwidth}\footnotesize \textit{Notes.} This table reports cumulative raw returns, the compound product of daily returns over each event window. All series are constructed from CRSP daily returns. The market return is value-weighted using previous-day market capitalization and reproduces the CRSP value-weighted index. Firms are classified by the gold-clause exposure measure $d_j$ defined in equation~(\ref{eq:d}): the two equal-weighted portfolios contain firms with $d_j>0$ (177 firms) and $d_j=0$ (376 firms), respectively, and the exposure-weighted portfolio weights firms with $d_j>0$ by $w_j = d_j/\sum_k d_k$. Portfolios are rebalanced daily over firms with a return on each day. For the Supreme Court decision (February~18), the benchmark is the close of February~16---the last trading day before the ruling, as the exchange was open on Saturdays. The modest decline during the oral arguments (January~8--10) understates the market anxiety triggered by the government's weak showing: press coverage in the following days drove sustained selling, with Liberty Bond prices reaching their highest level since 1917 by January~13 \citep{NYT1935}. Internet Appendix Section~\ref{sec:ia_other_events} examines the remaining legal and political events of the litigation period and reports the corresponding return differentials.\end{minipage}
\end{table}
